\begin{tikzpicture}[x=1cm,y=1cm,>=Latex,
 band/.style={draw=SPGrid,rounded corners=1.2mm,fill=SPFloor,
  minimum width=11.40cm,minimum height=2.48cm},
 card/.style={draw=black!23,rounded corners=.9mm,fill=white,
  minimum width=2.35cm,minimum height=1.18cm,align=center,font=\tiny},
 stage/.style={draw=black!23,rounded corners=.9mm,fill=white,
  minimum width=2.12cm,minimum height=1.18cm,align=center,font=\tiny},
 arrow/.style={->,thick,draw=VIUOrange},
 down/.style={->,thick,draw=black!40}]

\node[font=\scriptsize\bfseries,text=black!70] at (0,4.80)
 {MISMO MUNDO FÍSICO; TRES ESTADOS ESTRATÉGICOS};
\node[font=\tiny,text=black!55] at (0,4.49)
 {$a$ es ejecutable; $x$ y $(x,\lambda)$ pasan por $\mathcal R$; F-IV produce una secuencia $a_t$};

% ================================================================ F-II
\node[band] (fii) at (0,3.12) {};
\node[anchor=north west,font=\scriptsize\bfseries,text=SPBlue]
 at ($(fii.north west)+(.18,-.13)$) {F-II · INTENCIÓN CONTINUA $x_i\in\Delta_i$};
\node[card] (rep) at (-4.12,3.14)
 {\textbf{Replicator}\\[.6mm]
  $x_i=(.7,.3,0)$\\
  una componente en cero\\
  no se reactiva};
\node[card] (smith) at (-1.38,3.14)
 {\textbf{Smith}\\[.6mm]
  $k_1\rightleftarrows k_2$\\
  transfiere masa según\\
  diferencias de fitness};
\node[card] (bnn) at (1.38,3.14)
 {\textbf{BNN}\\[.6mm]
  $\widehat F_{ik}=[F_{ik}-\bar F_i]_+$\\
  entra masa donde el\\
  exceso es positivo};
\node[card] (logit) at (4.12,3.14)
 {\textbf{Logit}\\[.6mm]
  $x_i^+\propto\exp(F_i/\tau)$\\
  softmax difuso; $\tau$\\
  regula exploración};
\node[draw=VIUOrange,rounded corners=.8mm,fill=VIUOrange!4,
 minimum width=4.15cm,minimum height=.46cm,font=\tiny\bfseries] (closure) at (0,2.05)
 {$x^\star\quad\xrightarrow{\ \mathcal R\ }\quad a^\star$ entero};
\foreach \n in {rep,smith,bnn,logit}{\draw[down] (\n.south)--(closure.north);}
\node[font=\fontsize{5.3}{5.6}\selectfont,text=SPRed] at (0,1.53)
 {las cuatro dinámicas comparten $F$ y el cierre; el ranking puede cambiar después de $\mathcal R$};

% ================================================================ F-III
\node[band] (fiii) at (0,.00) {};
\node[anchor=north west,font=\scriptsize\bfseries,text=SPGreen!70!black]
 at ($(fiii.north west)+(.18,-.13)$) {F-III · CUOTAS COMO RESTRICCIÓN COMPARTIDA};
\node[stage,draw=VIUOrange!70] (scarcity) at (-4.18,.02)
 {\textbf{1 · escasez}\\[.6mm]
  carga $k_1$\\
  $Q_1/m_1=2/5$\\
  $g_1(x)>0$};
\node[stage,draw=SPRed!60] (price) at (-1.40,.02)
 {\textbf{2 · precio}\\[.6mm]
  $\lambda_1\uparrow$\\
  $c_i\lambda_1-d_{i1}$\\
  atrae capacidad};
\node[stage,draw=SPPurple!60] (consensus) at (1.40,.02)
 {\textbf{3 · consenso}\\[.6mm]
  $\lambda_{11}\leftrightarrow\lambda_{21}$\\
  $\leftrightarrow\lambda_{31}$\\
  copias vecinales};
\node[stage,draw=SPGreen!65] (physical) at (4.18,.02)
 {\textbf{4 · salida}\\[.6mm]
  $(x^\star,\lambda^\star)$\\
  $\mathcal R(x^\star)=a^\star$\\
  coalición atómica};
\draw[arrow] (scarcity.east)--(price.west);
\draw[arrow] (price.east)--(consensus.west);
\draw[arrow] (consensus.east)--(physical.west);
\node[font=\tiny,text=black!55] at (0,-1.00)
 {aceptación numérica: residual primal + dual + complementariedad + estación + consenso};

% ================================================================ F-IV
\node[band] (fiv) at (0,-3.12) {};
\node[anchor=north west,font=\scriptsize\bfseries,text=SPPurple]
 at ($(fiv.north west)+(.18,-.13)$) {F-IV · RECLUTAMIENTO BAJO EVENTOS};
\node[stage] (t0) at (-4.18,-3.10)
 {\textbf{$t_0$ · servicio}\\[.6mm]
  cargas activas\\
  robots disponibles\\
  coalición $a_t$};
\node[stage,draw=VIUOrange!65] (t1) at (-1.40,-3.10)
 {\textbf{$t_1$ · llegada}\\[.6mm]
  aparece $k_{\rm new}$\\
  cambia la demanda\\
  estado $s_{t+1}$};
\node[stage,draw=SPRed!65] (t2) at (1.40,-3.10)
 {\textbf{$t_2$ · fallo}\\[.6mm]
  $r_j$ deja de estar\\
  disponible\\
  se invalida parte de $a_t$};
\node[stage,draw=SPGreen!65] (t3) at (4.18,-3.10)
 {\textbf{$t_3$ · política}\\[.6mm]
  reactiva / inicialización previa\\
  oráculo escalar finito\\
  nueva asignación $a_{t+1}$};
\draw[arrow] (t0.east)--node[above,font=\tiny] {evento}(t1.west);
\draw[arrow] (t1.east)--node[above,font=\tiny] {$T(s,u,e)$}(t2.west);
\draw[arrow] (t2.east)--node[above,font=\tiny] {decisión}(t3.west);
\node[font=\tiny,text=SPRed] at (0,-4.14)
 {estudio temporal con eventos exógenos; sin garantía de potencial de Markov};

\end{tikzpicture}
